\documentclass[../cours_1re.tex]{subfiles}
\begin{document}
\setcounter{chapter}{20}
\chapter{Variables aléatoires et espérance}
\label{ch:21}
\begin{objectif}
Définir une variable aléatoire sur un univers fini, déterminer sa loi et démontrer la linéarité de l'espérance.
Prérequis : probabilités finies, sommes et fonctions.
\end{objectif}
\section{Une fonction dont l'entrée est aléatoire}
\begin{definition}
Une variable aléatoire réelle $X$ sur un univers fini $\Omega$ est une fonction
$X:\Omega\to\R$. Sa loi donne les valeurs distinctes $x_1,\ldots,x_r$
et les probabilités $p_i=P(X=x_i)$.
\end{definition}
L'événement $\{X=x_i\}$ est l'ensemble des issues $\omega$ telles que
$X(\omega)=x_i$. Ces événements partitionnent $\Omega$, donc $\sum_ip_i=1$.
$P(X\le a)$ s'obtient en additionnant les $p_i$ pour lesquels $x_i\le a$.
Il faut distinguer l'issue $\omega$, sa valeur $X(\omega)$ et la probabilité
de cette valeur.

\begin{exemple}
On lance deux dés équilibrés indépendants et $X$ est leur somme.
Les $36$ couples ordonnés sont équiprobables. Les effectifs des sommes
$2,3,\ldots,12$ sont $1,2,3,4,5,6,5,4,3,2,1$.
Donc $P(X=7)=6/36$ et $P(X\ge10)=(3+2+1)/36=1/6$.
Les onze sommes ne sont pas équiprobables.
\end{exemple}

\section{Espérance et interprétation}
\begin{definition}
$\displaystyle \mathbb E(X)=\sum_i x_iP(X=x_i)$.
C'est une moyenne pondérée des valeurs possibles.
\end{definition}
En regroupant les issues qui donnent la même valeur, on obtient également
$\mathbb E(X)=\sum_{\omega\in\Omega}X(\omega)p_\omega$.
Cette seconde forme facilite les démonstrations.
\begin{propriete}[linéarité]
Pour des variables $X,Y$ sur un même univers et des réels $a,b$ :
$\mathbb E(aX+bY)=a\mathbb E(X)+b\mathbb E(Y)$.
Aucune indépendance n'est nécessaire.
\end{propriete}
\begin{demonstration}
Développer la somme
$\sum_\omega(aX(\omega)+bY(\omega))p_\omega$ et distribuer les termes
donne exactement les deux sommes annoncées.
Une constante $c$ a pour espérance $c\sum_\omega p_\omega=c$.
\end{demonstration}
Si $X\le Y$ pour chaque issue de probabilité positive, alors
$\mathbb E(X)\le\mathbb E(Y)$, car la somme des termes
$(Y(\omega)-X(\omega))p_\omega$ est positive ou nulle.
L'espérance se situe donc entre le minimum et le maximum des valeurs possibles.
Elle n'est pas nécessairement une valeur prise : un dé a une espérance $3{,}5$.

\section{Variables indicatrices}
La variable $\mathbf1_A$ vaut $1$ sur $A$, $0$ ailleurs.
Sa loi donne immédiatement $\mathbb E(\mathbf1_A)=P(A)$.
Pour le nombre $X$ de succès parmi $n$ épreuves de probabilité $p$,
$X=\mathbf1_{S_1}+\cdots+\mathbf1_{S_n}$, donc
$\mathbb E(X)=np$. Cette espérance reste vraie même sans indépendance
si chaque succès a encore la probabilité $p$ ; c'est la loi binomiale,
et plus tard la formule simple de variance, qui exigent l'indépendance.
\begin{exemple}
Pour deux dés, la linéarité donne $\mathbb E(X)=3{,}5+3{,}5=7$,
sans additionner les onze termes de la loi de la somme.
\end{exemple}

\section{Jeux et décisions : préciser ce que l'on mesure}
Un jeu verse $20$ euros avec probabilité $1/10$ et $0$ sinon.
Si la mise est $m$, le gain net $G$ vaut $20-m$ ou $-m$ ;
$\mathbb E(G)=2-m$. Il est équitable au sens de l'espérance si $m=2$.
Cela ne garantit aucun résultat sur une partie ni une absence de risque.
Le mot « gain » doit préciser s'il inclut le remboursement de la mise.

\section{Transformer une variable}
Pour toute fonction $f$ définie sur les valeurs de $X$,
\[
\mathbb E(f(X))=\sum_i f(x_i)p_i.
\]
Pour le démontrer, sommer d'abord sur les issues puis regrouper suivant $X$.
En général $\mathbb E(f(X))\ne f(\mathbb E(X))$.
Si $X$ vaut $-1$ ou $1$ avec probabilités égales,
$\mathbb E(X)=0$, mais $\mathbb E(X^2)=1$.
L'égalité est garantie pour les fonctions affines par linéarité,
pas pour un carré ni une exponentielle.

\begin{methode}
Pour déterminer une loi : définir l'expérience et son univers ; définir la
variable ; lister ses valeurs distinctes ; calculer chaque probabilité ;
vérifier leur somme ; seulement ensuite calculer les moments demandés.
\end{methode}

\begin{finchapitre}
\section{Utilité et prolongements}
L'espérance convertit une expérience en une quantité moyenne, mais la
dispersion demande un second outil : la variance.

En licence, l'espérance sera une intégrale. Sa linéarité restera fondamentale,
mais les sommes infinies demanderont des conditions de convergence.
Les indicatrices relieront toujours les fonctions aux événements.
\begin{retenir}
Une variable aléatoire est une fonction. La linéarité de l'espérance ne
demande pas l'indépendance. Une moyenne théorique n'est pas une garantie individuelle.
\end{retenir}
\end{finchapitre}
\end{document}
